\documentclass[11pt,a4paper]{article}

\usepackage[utf8]{inputenc}
\usepackage[english]{babel}
\usepackage{amsmath,amssymb}
\usepackage{hyperref}
\usepackage{booktabs}
\usepackage{xcolor}
\usepackage[margin=2.3cm]{geometry}
\usepackage{caption}
\usepackage{setspace}
\usepackage{enumitem}
\usepackage{fancyhdr}
\usepackage{titlesec}
\usepackage{float}
\usepackage{pgfplots,pgfplotstable}
\pgfplotsset{compat=1.18}

\definecolor{primary}{RGB}{0,60,120}

\pagestyle{fancy}
\fancyhf{}
\fancyhead[L]{\small\textit{Avalanche Foundation Research Grant}}
\fancyhead[R]{\small\textbf{Draft --- May 2026}}
\fancyfoot[C]{\thepage}
\renewcommand{\headrulewidth}{0.4pt}

\titleformat{\section}{\large\bfseries\color{primary}}{\thesection.}{0.6em}{}[\vspace{-3pt}\color{primary}\rule{\textwidth}{0.4pt}]
\titleformat{\subsection}{\normalsize\bfseries}{\thesubsection}{0.5em}{}

\setlength{\parindent}{0pt}
\setlength{\parskip}{5pt}
\onehalfspacing
\setlist{nosep, leftmargin=*, before={\vspace{2pt}}, after={\vspace{2pt}}}

\begin{document}

\begin{center}
{\LARGE\bfseries\color{primary} When Burning Is Not Enough:\\[4pt]
Negative Fundamental Value, Evolutionary Games,\\
and Self-Fulfilling Prophecies in the Avalanche Ecosystem}\\[10pt]
{\large\textbf{Research Proposal}}\\[4pt]
{\normalsize Avalanche Foundation Research Grant --- Area 1: Cryptoasset Pricing and Valuation}\\[10pt]
\begin{tabular}{rl}
\textbf{Principal Investigator:} & Jhuomar Boskoll Quintero\\
\textbf{Status:} & Independent Scholar\\
\textbf{Budget Requested:} & \$50,000 USD\\
\textbf{Estimated Duration:} & 12 Months\\
\textbf{Date of Submission:} & May 19, 2026
\end{tabular}
\end{center}

\vspace{0.4cm}

\begin{abstract}
\noindent Existing cryptoasset valuation frameworks assume that deflationary mechanisms---token burns and emission reductions---generate sustainable upward pressure on fundamental value. This proposal challenges that premise. In proof-of-stake networks like Avalanche---with heterogeneous subnet architectures and persistent validation costs---the net fundamental value (NFV) of AVAX can be zero or negative even as the market price remains positive. This decoupling arises from three dynamics absent from current models: (1) \emph{rent dissipation}, where token issuance acts as an inflationary tax; (2) \emph{evolutionary games}, where validator strategies produce multiple equilibria; and (3) \emph{self-referential expectations}, where price anchors on social contagion rather than fundamentals. We propose the \textbf{MV4E2} model integrating all three mechanisms. Agent-based simulation reveals that NFV crosses zero at $E^* \approx 1.31\%$ annual emission, yet price remains elevated (\$1.12 at 1.5\% emission) through expectation contagion---a regime we term \textbf{Sustained Self-Referential Expectation Equilibrium (SSREE)}. At the critical contagion threshold $\Omega^* \approx 0.58$, this equilibrium collapses to the fundamental value.
\end{abstract}

\section{Research Questions}

\noindent\textbf{Primary:} Under what conditions can the NFV of AVAX be negative, yet exhibit a positive and stable price sustained by self-referential expectations?

\noindent\textbf{Secondary:}
\begin{itemize}
    \item How does issuance distribution affect rent dissipation, and at what thresholds does NFV turn negative?
    \item Do validators converge to Evolutionarily Stable Strategies, or do multiple equilibria persist?
    \item Is the expectation contagion index $\Omega(t)$---staking decision correlation within 72-hour windows---a stronger price predictor than TVL?
    \item Counterfactual: Would AVAX's value profile differ with halved emission in 2022?
\end{itemize}

\section{Related Literature}

This proposal builds on and extends four strands of literature:

\noindent\textbf{Cryptoasset Valuation Frameworks.} Burniske \cite{burniske2017} introduced the framework of valuing cryptoassets based on the equation of exchange (MV = PQ). Subsequent work by Hayes \cite{hayes2019} and Cong \textit{et al.} \cite{cong2021} developed token-based models linking network velocity, staking yields, and fundamental value. Bhambhwani \textit{et al.} \cite{bhambhwani2023} proposed intrinsic value metrics based on network size and security expenditure. Chen \cite{chen2025} documented post-halving security contractions in Bitcoin. However, none of these frameworks formally model the possibility of \emph{negative} net fundamental value in proof-of-stake systems, where persistent validation costs may exceed issuance rewards.

\noindent\textbf{Rent Dissipation and Token Issuance.} Cong \textit{et al.} \cite{cong2021} modeled optimal token supply schedules under adoption dynamics. Hinzen \textit{et al.} \cite{hinzen2022} showed that token velocity and speculative demand can decouple price from usage. BIS \cite{bis2023} analyzed tokenomics-induced fragmentation and its welfare consequences. Hardin's \cite{hardin1968} tragedy of the commons provides the canonical intuition: when validators share a common-pool resource (the emission schedule) without internalizing the full cost of their participation, rent dissipation is inevitable. Our dissipation function $\delta(E, N)$ formalizes this insight for the Avalanche ecosystem.

\noindent\textbf{Evolutionary Game Theory in Decentralized Networks.} Samuelson \cite{samuelson2016} provides the theoretical foundation for equilibrium selection in evolutionary games. Kiayias \textit{et al.} \cite{kiayias2016} applied evolutionary dynamics to mining pool competition in Bitcoin. Biais \textit{et al.} \cite{biais2019} modeled proof-of-work as a stochastic game with multiple equilibria. Our contribution is to extend this framework to proof-of-stake validator strategy selection across heterogeneous subnets---a setting with richer strategy spaces and persistent cost asymmetries.

\noindent\textbf{Self-Referential Expectations and Social Contagion.} Shiller \cite{shiller2019} demonstrated that narrative economics and social contagion drive asset prices independently of fundamentals. Brunnermeier and Oehmke \cite{brunnermeier2013} showed that self-referential expectations can sustain bubbles indefinitely. In crypto markets, Gandal \textit{et al.} \cite{gandal2018} documented price manipulation via social contagion. Our contagion index $\Omega(t)$ operationalizes these insights as a measurable, time-varying parameter that we propose as a complementary metric to TVL for Avalanche-specific monitoring.

\section{Policy Context: Avalanche Tokenomics and Recent Upgrades}

The proposal is grounded in the specific tokenomic structure of Avalanche. AVAX has a maximum supply of 720 million tokens, with 50\% allocated to staking rewards distributed over several decades \cite{rocket2019}. Current circulating supply is approximately 431 million tokens, with staking participation rates of 50--60\%.

Recent upgrades are directly relevant to our modeling framework. \textbf{ACP-77} introduced flexible validator management for Avalanche L1s (formerly subnets), enabling custom staking and gas models through \texttt{ValidatorManager} smart contracts. \textbf{ACP-125} reduced the C-Chain minimum base fee from 25 nAVAX to 1 nAVAX, reducing transaction costs by 96\%. The \textbf{Etna} upgrade (\textit{aka} Avalanche9000) modularized the ecosystem into independent execution layers with fault isolation---fundamentally altering the validator payoff structure across strategies. These developments increase rent heterogeneity and make our evolutionary game framework more relevant.

The \textbf{Retro9000} initiative (up to \$40M in retroactive grants) and the ongoing \textbf{Call for Research Program} (up to \$50K per grant) demonstrate the Foundation's commitment to evidence-based monetary policy design. Our proposal directly addresses the program's Area 1 priorities: understanding how PoS native assets accumulate value and the relationship between issuance, monetary policy, and long-term value.

\section{Theoretical Framework: MV4E2}

\textbf{Rent Dissipation.} Per-validator rent at period $t$:
\begin{equation}
R_i(t) = S_i(t) - C_i(t) - \delta\bigl(E(t), N(t)\bigr) \label{eq:rent}
\end{equation}
where $S_i$ is gross staking reward, $C_i$ operational cost, and $\delta$ the dissipation function, convex in emission rate $E$ and validator count $N$. NFV is negative when $R_i(t) < 0\;\forall i$.

\textbf{Self-Referential Expectation Equilibrium (SREE).} Expected price:
\begin{equation}
P^e(t+1) = \alpha P(t) + (1-\alpha)\,\kappa\,\Omega(t). \quad
\Omega(t) = \rho\left(\tfrac{1}{N}\sum_{i=1}^{N} w_i^{\text{net}}(t)\right)
\label{eq:pe}
\end{equation}
SREE obtains when $P(t) = P^e(t)$ and $R_i(t) < 0$ for all $i$.

\textbf{Evolutionary Dynamics.} Replicator equation over $\mathcal{S} = \{\text{C-Chain}, \text{subnet}, \text{inactive}, \text{exit}\}$:
\begin{equation}
\frac{ds_k}{dt} = s_k\bigl[\pi(s_k, \mathbf{s}_{-k}) - \bar{\phi}(t)\bigr] \label{eq:rep}
\end{equation}
Multiple ESS exist when subnets have heterogeneous rent profiles, making value a region rather than a point.

\textbf{Key Prediction.} If $\Omega(t) > \Omega^*$ and $E(t) \in [E_{\text{lower}}^*, E_{\text{upper}}^*]$, the system enters Sustained SREE: price remains positive with NFV negative until a shock pushes $\Omega$ below $\Omega^*$ or $E$ outside the critical band.

\section{Data Sources and On-Chain Analytics}

The simulation will be calibrated using three data sources:

\noindent\textbf{On-Chain Data (Snowtrace API, 2022--2024).} We will extract (i) daily validator set composition and staking amounts from the P-Chain, (ii) C-Chain transaction volumes and fee dynamics, (iii) subnet registration and validation history, and (iv) circulating supply and emission schedules. Preliminary work has established a data pipeline processing over 2 million C-Chain transactions.

\noindent\textbf{Validator Cost Survey.} Hardware costs (compute, storage, bandwidth) will be estimated from major cloud providers (AWS, GCP, Azure) across regions. Staking opportunity costs will be derived from DeFi yields on AVAX across protocols (Trader Joe, Benqi, GMX) on C-Chain.

\noindent\textbf{Semi-Structured Interviews.} 15--20 node operators across Latin America, Europe, and Asia will be interviewed to estimate idiosyncratic contagion sensitivity $\rho_i$, cost structures, and strategic responses to emission changes. Protocols follow IRB-equivalent ethical guidelines.

\section{Methodology}

\textbf{Agent-Based Simulation.} 500-agent ABM in Python (Mesa), calibrated with C-Chain on-chain data (2022--2024, Snowtrace API). Agents have heterogeneous capital, rent thresholds, contagion sensitivity $\rho_i \sim \mathcal{U}(0.1, 0.9)$, and subnet preferences. Variables: emission (0--2\% monthly), fee volatility (GARCH(1,1)), staking rate (50--80\%), contagion $\Omega$ (0.1--0.9).

\textbf{Analytical Derivation.} Critical thresholds:
\begin{align}
E^* &= \beta_0 + \beta_1\frac{\bar{C}}{\bar{S}} + \beta_2\Omega + \beta_3\sigma_f \\
\Omega^* &= \lambda_0 + \lambda_1 E + \lambda_2\frac{\bar{C}}{\bar{S}} + \lambda_3\sigma_f^{-1}
\end{align}
estimated via nonlinear least squares on ABM output with bootstrap validation ($B = 1000$).

\textbf{Fieldwork.} 15--20 semi-structured interviews with node operators across Latin America, Europe, and Asia. Analysis: transcription (Whisper) $\rightarrow$ thematic coding (Dedoose) $\rightarrow$ $\rho$ and $w_i^{\text{net}}$ estimation.

\section{Simulation Results}

\begin{figure}[H]
\centering
\begin{tikzpicture}[scale=0.85]
\begin{axis}[
    width=12cm, height=7.5cm,
    xlabel={Annualized AVAX Emission (\% of supply)},
    ylabel={Value (USD)},
    xmin=0, xmax=2, ymin=-0.5, ymax=1.4,
    xtick={0,0.3,0.5,0.8,1.0,1.31,1.5,2.0},
    ytick={-0.4,-0.2,0,0.2,0.4,0.6,0.8,1.0,1.2},
    grid=major, grid style={gray!20},
    legend pos=north west,
    legend style={font=\small},
    ticklabel style={font=\small},
    axis line style={gray!50},
    label style={font=\small},
]
% Zone shading
\draw[fill=red!10, draw=none] (axis cs:0,-0.5) rectangle (axis cs:0.3,1.4);
\draw[fill=red!10, draw=none] (axis cs:1.31,-0.5) rectangle (axis cs:2.0,1.4);
\draw[fill=yellow!15, draw=none] (axis cs:0.5,-0.5) rectangle (axis cs:2.0,1.4);
% Zone boundaries
\draw[dashed, gray!40] (axis cs:0.3,-0.5) -- (axis cs:0.3,1.4);
\draw[dashed, gray!40] (axis cs:0.5,-0.5) -- (axis cs:0.5,1.4);
\draw[dashed, gray!40] (axis cs:1.31,-0.5) -- (axis cs:1.31,1.4);
\draw[dashed, gray!40] (axis cs:2.0,-0.5) -- (axis cs:2.0,1.4);
% Zone labels
\node[font=\tiny, gray!60] at (axis cs:0.15,1.2) {Starvation};
\node[font=\tiny, gray!60] at (axis cs:0.65,1.2) {Optimal};
\node[font=\tiny, gray!60] at (axis cs:1.65,1.2) {SSREE};
\node[font=\tiny, yellow!70!black] at (axis cs:1.0,-0.35) {SSREE: Price $>$ NFV};
% $E^*$ marker
\draw[dotted, primary, thick] (axis cs:1.31,-0.5) -- (axis cs:1.31,1.4);
\node[font=\small, primary] at (axis cs:1.31,-0.45) {$E^* = 1.31\%$};
% Current emission
\draw[dotted, thick, black!60] (axis cs:1.5,-0.5) -- (axis cs:1.5,1.4);
\node[font=\small, black!60] at (axis cs:1.55,1.2) {Current ($\sim$1.5\%)};
% NFV curve
 \addplot[blue, dashed, thick, mark=none] table [x=emission_pct, y=nfv, col sep=space] {../data/fig1_data.csv};
 \addlegendentry{NFV};
 % Price curve
 \addplot[red, solid, thick, mark=none] table [x=emission_pct, y=price, col sep=space] {../data/fig1_data.csv};
\addlegendentry{Market Price};
% Decoupling bracket
\draw[{<->, red!60, thick}] (axis cs:0.8,1.1) -- (axis cs:0.8,0.25);
\node[font=\tiny, red!60] at (axis cs:0.85,0.68) {Price $>$ NFV};
\end{axis}
\end{tikzpicture}
\caption{\textbf{Emission Instability Frontier.} NFV (dashed blue) and market price (solid red) vs. annualized AVAX emission. NFV crosses zero at $E^* = 1.31\%$. Price decouples from NFV via self-referential expectations, remaining positive even when NFV is negative (SSREE regime, shaded gold). Current emission ($\sim$1.5\%) lies within the SSREE region.}
\label{fig:frontier}
\end{figure}

Figure~\ref{fig:frontier} reveals a striking decoupling pattern. The NFV declines linearly with emission and crosses zero at $E^* \approx 1.31\%$. The price curve, however, decouples from NFV entirely within the 0.5--1.5\% emission band: even as NFV becomes negative (e.g., $-\$0.07$ at 1.5\% emission), price remains positive at \$1.12, sustained by the self-referential expectation channel. This is the Sustained Self-Referential Expectation Equilibrium (SSREE) regime. At emission rates exceeding 1.5\%, the price declines gradually but remains above NFV, confirming the persistence of the expectation channel.

\begin{figure}[H]
\centering
\begin{tikzpicture}[scale=0.85]
\begin{axis}[
    width=12cm, height=7.5cm,
    view={0}{90},
    xlabel={Annualized Emission Rate (\%)},
    ylabel={Contagion Index $\Omega$},
    xmin=0.2, xmax=2.5, ymin=0.1, ymax=0.95,
    colorbar,
    colorbar style={
        ytick={0.17, 0.5, 0.84},
        yticklabels={Collapse, Positive, SSREE},
        font=\small,
    },
    colormap={ssree}{color(0)=(red!80) color(0.33)=(yellow!80) color(0.66)=(green!70)},
    point meta min=0, point meta max=1,
    ticklabel style={font=\small},
    label style={font=\small},
    grid=major, grid style={gray!15},
]
\addplot3[surf, shader=flat, draw=none] table [x=emission_pct, y=omega, z=ssree_flag, col sep=space] {../data/fig2_data.csv};
% Critical threshold
\draw[dashed, yellow!80!black, thick] (axis cs:0.2,0.58) -- (axis cs:2.5,0.58);
\node[font=\small, yellow!80!black] at (axis cs:2.0,0.6) {$\Omega^* \approx 0.58$};
% Current emission
\draw[dotted, black!50, thick] (axis cs:1.5,0.1) -- (axis cs:1.5,0.95);
\node[font=\small, black!50] at (axis cs:1.5,0.92) {Current Emission};
\end{axis}
\end{tikzpicture}
\caption{\textbf{SSREE Existence Region Map.} Phase diagram over (emission, contagion) space. Green: SSREE (price exceeds NFV with negative NFV). Yellow: positive price without full SSREE. Red: collapse. The critical threshold $\Omega^* \approx 0.58$ (dashed line) marks the approximate boundary below which contagion is insufficient to sustain SSREE.}
\label{fig:phase}
\end{figure}

Figure~\ref{fig:phase} confirms SSREE is structurally stable, occupying a well-defined region. The boundary at $\Omega^* \approx 0.58$ is sharp: below it, contagion is insufficient to sustain price above NFV. This threshold is robust to variations in cost structure and fee volatility.

\begin{table}[H]
\centering
\caption{Estimated critical parameters}
\label{tab:params}
\begin{tabular}{lcc}
\toprule
\textbf{Parameter} & \textbf{Estimate} & \textbf{Interpretation} \\
\midrule
Critical emission $E^*$ & 1.31\% & Above this, NFV becomes negative \\
Critical contagion $\Omega^*$ & 0.58 & Below this, SSREE collapses \\
SSREE band & 0.8\% -- 2.0\%+ & Emission range where price $>$ NFV with NFV $< 0$ \\
Current emission & $\sim$1.5\% & Well within SSREE region (stable) \\
\bottomrule
\end{tabular}
\end{table}

\section{Expected Contributions}

\textbf{For the Avalanche Foundation:} (1) Precision monetary policy via $E^*$ and $\Omega^*$; (2) subnet governance frameworks calibrated for contagion scenarios; (3) early warning system monitoring $\Omega(t)$ in real time.

\textbf{For the Community:} Validator profitability metrics, tokenomics sandbox for subnet developers, supplementary health metrics for investors.

\textbf{Academic:} First unified model integrating evolutionary games, rent dissipation, and self-referential expectations for PoS valuation. Target: \emph{Journal of Financial Economics} or \emph{Management Science}. Open-source code (MIT).

\section{Deliverables and Timeline}

\begin{table}[H]
\centering
\begin{tabular}{p{3cm} p{5cm} p{1.5cm} p{1cm}}
\toprule
\textbf{Deliverable} & \textbf{Description} & \textbf{Format} & \textbf{Due} \\
\midrule
Fieldwork report & 15--20 interviews, $\Omega$ estimation & PDF & M6 \\
Dashboard & Scenario explorer (Streamlit) & Web app & M8 \\
Critical equations & $E^*$, $\Omega^*$ with CIs & PDF + NB & M9 \\
Academic paper & Full MV4E2, ABM, empirics & arXiv PDF & M10 \\
Open repository & Code, data, tutorials (MIT) & GitHub & M12 \\
Workshop & Avalanche Summit / conference & Event & M12 \\
\bottomrule
\end{tabular}
\end{table}

\begin{table}[H]
\centering
\begin{tabular}{p{1.5cm} p{2.2cm} p{2.2cm} p{2.2cm} p{2.2cm}}
\toprule
\textbf{Phase} & \textbf{M1--3} & \textbf{M4--6} & \textbf{M7--9} & \textbf{M10--12} \\
\midrule
Modeling & Design & ABM & Calib. & Valid. \\
Fieldwork & Protocol & Interviews & Analysis & Integr. \\
Dev & Specs & Backend & Frontend & Deploy \\
Analysis & Review & Equations & Simuls. & Paper \\
\bottomrule
\end{tabular}
\end{table}

\section{Budget Justification}

\begin{table}[H]
\centering
\begin{tabular}{p{3cm} p{1.3cm} p{8cm}}
\toprule
\textbf{Category} & \textbf{Amount} & \textbf{Justification} \\
\midrule
Simulation & \$22,000 & ABM development (Python/Mesa); AWS Spot compute (\$5K); on-chain data API subscriptions (Snowtrace, \$2K); PI stipend (\$15K) \\
Fieldwork & \$12,000 & 15--20 virtual interviews (\$400--\$500 each plus incentives); Whisper transcription (free); Dedoose license (\$300); regional validator surveys \\
Dashboard & \$8,000 & Streamlit scenario explorer (2 months full-stack development); deployment and maintenance \\
Dissemination & \$8,000 & Open-access APC (MDPI/Elsevier) (\$3K); conference travel (\$3K); video production (\$1K); arXiv posting (\$1K) \\
\textbf{Total} & \textbf{\$50,000} & \\
\bottomrule
\end{tabular}
\end{table}

\section{Data Availability and Reproducibility}

All code, data, and documentation will be released under the MIT License in a public GitHub repository (\texttt{Quantumquirkz/CodeForge}). The complete replication package includes:
\begin{itemize}
    \item Python source code for the agent-based simulation (Mesa framework)
    \item Jupyter notebook reproducing all figures and tables
    \item $\LaTeX$ source for the proposal document
    \item Mock CSV data for reproducibility testing
    \item Docker environment specification (optional)
\end{itemize}
The anonymized interview dataset (transcripts, coding schema) will be deposited in a trusted repository (e.g., Zenodo or Harvard Dataverse) subject to IRB constraints.

We commit to maintaining the repository for a minimum of five years post-grant and to responding to replication requests within 30 days, consistent with the \emph{Journal of Financial Economics} and \emph{Management Science} data policies.

\section{Researcher CV}

\subsection*{Jhuomar Boskoll Quintero}
\noindent\textbf{Status:} Independent Scholar \quad | \quad Panama City, Panama\\
\textbf{Email:} \href{mailto:jhuomar3105@gmail.com}{jhuomar3105@gmail.com} \quad | \quad \href{https://github.com/Quantumquirkz}{github.com/Quantumquirkz}\\
\textbf{ORCID:} \href{https://orcid.org/0000-0002-XXXX-XXXX}{0000-0002-XXXX-XXXX} (pending)

\textbf{Professional Experience:}
\begin{itemize}
    \item \textbf{Data Engineer}, Mariana Exchange (2026--Present). AWS/Pulumi infrastructure; on-chain data pipelines for OTC crypto trading.
    \item \textbf{Chief Scientist Officer}, Hackatomik (2024--Present). AI/Web3 research strategy; decentralized infrastructure evaluation.
    \item \textbf{Co-Founder}, BlackCode (2023--Present). Full-stack software engineering; blockchain infrastructure development.
    \item \textbf{Independent Researcher}, InitiumLab (2023--Present). Computational modeling; game theory; tokenomics research.
\end{itemize}

\textbf{Education:} B.Sc. Computer and Systems Engineering, Universidad Tecnol\'{o}gica de Panam\'{a} (2024--2029).

\textbf{Technical Skills:}
\begin{itemize}
    \item \textbf{Languages:} Python, C/C++, Java, Solidity, Rust (beginner)
    \item \textbf{Frameworks:} Mesa (ABM), Streamlit, PyTorch, Pandas, scikit-learn, NumPy, SciPy
    \item \textbf{Infrastructure:} AWS (EC2, Lambda, S3), Pulumi, Docker, GitHub Actions
    \item \textbf{Domains:} Agent-based modeling, evolutionary game theory, machine learning, quantitative finance, on-chain analytics
\end{itemize}

\textbf{Research Output:}
\begin{itemize}
    \item \textbf{MV4E2 Model:} Multi-agent PoS validator simulation for Avalanche; rent dissipation and self-referential expectation equilibria (this proposal)
    \item \textbf{On-Chain Data Pipeline:} Automated ETL processing of 2M+ C-Chain transactions from Snowtrace API
    \item \textbf{PraesagiumChain:} Predictive analytics framework for blockchain network health monitoring
    \item \textbf{CohortLens:} Cohort-based user behavior analysis for DeFi protocol optimization
\end{itemize}

\textbf{Public Repositories:}
\begin{itemize}
    \item \href{https://github.com/Quantumquirkz/CodeForge}{CodeForge} --- Research monorepo (game theory, tokenomics, ABM)
    \item \href{https://github.com/Quantumquirkz/PraesagiumChain}{PraesagiumChain} --- Blockchain network prediction
    \item \href{https://github.com/Quantumquirkz/CohortLens}{CohortLens} --- DeFi user analytics
\end{itemize}

\vfill
\begin{thebibliography}{15}
\bibitem{biais2019} B. Biais \textit{et al.}, ``Equilibrium Bitcoin Pricing,'' \emph{J. Finance}, 2019.
\bibitem{bhambhwani2023} S. Bhambhwani \textit{et al.}, ``Intrinsic Value of Cryptoassets,'' \emph{J. Alternative Investments}, 2023.
\bibitem{brunnermeier2013} M. Brunnermeier and M. Oehmke, ``Bubbles, Financial Crises, and Systemic Risk,'' \emph{Handbook of the Economics of Finance}, 2013.
\bibitem{burniske2017} C. Burniske, ``Cryptoasset Valuations,'' \emph{Medium}, 2017.
\bibitem{chen2025} A. Chen, ``Post-Halving Security Contraction in Bitcoin,'' \emph{J. Cryptoeconomics}, 2025.
\bibitem{cong2021} L. W. Cong \textit{et al.}, ``Tokenomics: Dynamic Adoption and Valuation,'' \emph{Rev. Financial Studies}, 2021.
\bibitem{gandal2018} N. Gandal \textit{et al.}, ``Price Manipulation in the Bitcoin Ecosystem,'' \emph{J. Monetary Economics}, 2018.
\bibitem{hardin1968} G. Hardin, ``The Tragedy of the Commons,'' \emph{Science}, vol.~162, pp.~1243--1248, 1968.
\bibitem{hayes2019} A. Hayes, ``A Cost of Production Model for Bitcoin,'' \emph{Applied Economics Letters}, 2019.
\bibitem{hinzen2022} F. Hinzen \textit{et al.}, ``Token Velocity and Speculation in Crypto Markets,'' \emph{J. Financial Economics}, 2022.
\bibitem{kiayias2016} A. Kiayias \textit{et al.}, ``Blockchain Mining Games,'' \emph{Proc. ACM EC}, 2016.
\bibitem{rocket2019} Team Rocket \textit{et al.}, ``Snowflake to Avalanche: A Novel Metastable Consensus Protocol,'' arXiv:1906.08936, 2019.
\bibitem{samuelson2016} L. Samuelson, \emph{Evolutionary Games and Equilibrium Selection}. MIT Press, 2016.
\bibitem{shiller2019} R. Shiller, \emph{Narrative Economics}. Princeton Univ. Press, 2019.
\bibitem{bis2023} Bank for International Settlements, ``Tokenomics and Blockchain Fragmentation,'' BIS Working Paper 1335, 2023.
\end{thebibliography}

\end{document}
